\documentclass[11pt]{article}
\usepackage{mathblatt}
\newcommand{\bereich}[1]{\par\addvspace{14pt}\noindent{\large\bfseries #1}\par\nobreak\vspace{4pt}}
\newcommand{\bl}[1]{\par\noindent #1\par\vspace{3pt}}
\newcommand{\blrand}[1]{\par\noindent{\color{mbgrau}#1}\par\vspace{3pt}}
\newcommand{\blhier}[1]{\par\noindent\llap{$\blacktriangleright$\hspace{0.5em}}\textbf{#1}\par\vspace{3pt}}
\newcommand{\stufe}[1]{\par\noindent\hspace*{1.5em}{\small #1}\par\vspace{2pt}}
\newcommand{\pkt}{\textperiodcentered{}}
\begin{document}

\pfheftstil
\pflaufkopf[P10 \pkt{} Prüfung]{Wahrscheinlichk.}{PS50-5}\pfseitevon
\pfheftkopf{Wahrscheinlichk.}{}
\blhier{Wahrscheinlichkeit}
\stufe{Ergebnisse aufzählen}
\stufe{Wahrscheinlichkeit gesucht}
\stufe{Baum ergänzen}
\stufe{Pfadregeln anwenden}
\stufe{ohne Zurücklegen rechnen}
\stufe{Zufallsgerät entwerfen}

\par\addvspace{14pt}\noindent Vorher: Daten\par
\par\noindent Weiter: Zahlen\par
\end{document}